\section*{\centering RÉSUMÉ}
\addcontentsline{toc}{section}{RÉSUMÉ}

Les robots de service mobiles sont généralement livrés avec un écosystème logiciel fermé :
application propriétaire, absence de documentation et de protocole de communication publié.
Ce constat limite toute personnalisation par l'utilisateur final. Le projet \textsc{Cybel}, mené
dans le cadre d'un stage à HESTIM, vise à concevoir une plateforme de commande et
d'interaction indépendante pour un robot de réception mobile CIOT TY1251D-03195, composé d'un
châssis de navigation sous ROS et d'un \emph{« upper body »} Android.

La démarche adoptée repose sur une rétro-ingénierie incrémentale et non destructive du protocole
de communication interne du robot : balayage des services réseau exposés, introspection des
topics et services ROS via \texttt{rosbridge}/\texttt{rosapi}, et vérification systématique de
l'effectivité de chaque commande avant de la considérer comme fonctionnelle. Sur cette base, une
architecture en trois couches a été conçue~:
\begin{itemize}
    \item un SDK Python proposant une implémentation simulée (\emph{mock}) et une implémentation
    réelle interchangeables ;
    \item une API FastAPI (REST~+~WebSocket) ;
    \item une interface web Vite/TypeScript opérateur ;
\end{itemize}
complétée par une interface visiteur kiosque (\texttt{frontend-kiosk/}) et deux applications
Android natives légères (\texttt{CybelTTSBridge}, \texttt{CybelVisitorKiosk}).

À ce stade du stage, la connectivité avec le robot est établie, le protocole de télémétrie, de
commande de vitesse et de navigation autonome a été reconstruit et intégré, la synthèse vocale
(TTS) fonctionne via une application Android dédiée déclenchée par \texttt{am broadcast}, une
interface opérateur complète (carte SLAM, LiDAR, visiteurs détectés, gestion du parcours de
visite, arrêt d'urgence pendant une visite) est opérationnelle, et un déploiement embarqué sur
Termux (\emph{backend lite} + kiosque) est validé sur la tablette. L'interface visiteur a été
recentrée sur une visite guidée autonome du laboratoire : huit arrêts synchronisés depuis
\texttt{data/knowledgeV2-lab.json} vers \texttt{data/lab\_tour.json}, navigation par coordonnées
(\texttt{/navi\_goal}) et annonces vocales séquentielles.

Ce rapport présente le contexte, la problématique, l'état de l'art, la conception, la
méthodologie, l'implémentation réalisée, les résultats obtenus, ainsi qu'une analyse critique des
limites et perspectives du projet.

\bigskip
\noindent\textbf{Mots-clés :} robotique de service, rétro-ingénierie de protocole, ROS,
rosbridge, FastAPI, interface homme-robot, navigation autonome, synthèse vocale, Termux, WebView
Android.